

\documentclass[11pt,a4paper]{article}
\usepackage[hyperref]{emnlp2020}
\usepackage{times}
\usepackage{latexsym}
\renewcommand{\UrlFont}{\ttfamily\small}

\usepackage{microtype}

\aclfinalcopy 



\usepackage{subcaption}
\usepackage{tabularx}
\usepackage{booktabs}
\renewcommand{\UrlFont}{\ttfamily\small}
\usepackage{graphicx}
\usepackage{enumitem}
\usepackage{csvsimple}
\usepackage{dsfont}
\usepackage{pifont}
\usepackage{amsthm}
\usepackage[ruled,vlined]{algorithm2e}
\usepackage{xspace}
\usepackage{amsmath}
\usepackage{arydshln}
\usepackage{multirow}
\usepackage{colortbl}
\usepackage{float}
\usepackage{multicol}
\usepackage{array}
\usepackage{amssymb}




\setlist{nolistsep}


\newcommand\BibTeX{B\textsc{ib}\TeX}

\newcommand{\ours}{\textsc{SciFact}\xspace}
\newcommand{\oursClosed}{$\ours_{\mathrm{closed}}$\xspace}
\newcommand{\oursOpen}{$\ours_{\mathrm{open}}$\xspace}

\newcommand{\sysname}{\textsc{VeriSci}\xspace}


\newcommand{\fever}{\textsc{Fever}\xspace}
\newcommand{\eraser}{\textsc{ERASER}\xspace}
\newcommand{\snopes}{UKP Snopes\xspace}

\newcommand\dw[1]{{\color{blue}\{\textit{#1}\}$_{DW}$}}
\newcommand\ac[1]{{\color{teal}\{\textit{#1}\}$_{AC}$}}
\newcommand\kyle[1]{{\color{orange}\{\textit{#1}\}$_{kyle}$}}
\newcommand\klcomment[1]{{\color{orange}\{\textit{#1}\}$_{kyle}$}}
\newcommand\lucy[1]{{\color{purple}\{\textit{#1}\}$_{lucy}$}}
\newcommand{\hanna}[1]{\textcolor{cyan}{[#1]}}
\newcommand\peter[1]{{\color{purple}\{\textit{#1}\}$_{peter}$}}


\newcommand\todoit[1]{{\color{red}\{TODO: \textit{#1}\}}\xspace}
\newcommand\todo{{\color{red}{TODO}}\xspace}
\newcommand\needcite{{\color{red}{CITE}}\xspace}

\newcommand{\supported}{\textsc{Supported}\xspace}
\newcommand{\refuted}{\textsc{Refuted}\xspace}

\newcommand{\supports}{\textsc{Supports}\xspace}
\newcommand{\refutes}{\textsc{Refutes}\xspace}
\newcommand{\notenough}{\textsc{NoInfo}\xspace}

\newcommand{\psupports}{\textsc{PartiallySupports}\xspace}
\newcommand{\prefutes}{\textsc{PartiallyRefutes}\xspace}
\newcommand{\ppartial}{\textsc{Partial}\xspace}

\newcommand{\support}{\textsc{Support}\xspace}
\newcommand{\refute}{\textsc{Refute}\xspace}



\newcommand{\entailment}{entailment\xspace}
\newcommand{\Entailment}{Entailment\xspace}

\newcommand{\webpage}{\url{www.placeholder.com}}

\newcommand{\numclaims}{1,409\xspace}
\newcommand{\numpapers}{5,183\xspace}
\newcommand{\numevidence}{486\xspace}
\newcommand{\numcitednoevidence}{255\xspace}

\newcommand{\labelkappa}{0.75\xspace} \newcommand{\labelhumanacc}{0.83\xspace} \newcommand{\rationalekappaagree}{0.71\xspace} \newcommand{\rationalekappaoverall}{0.63\xspace} \newcommand{\rationalekappabackground}{0.45\xspace} \newcommand{\foneagree}{0.75\xspace}              \newcommand{\foneoverall}{0.67\xspace}           \newcommand{\nagreement}{232\xspace}              

\newcommand{\labelkappafine}{0.48\xspace} \newcommand{\nagreementfine}{224\xspace}              


\newcommand{\bestpipeline}{0.47\xspace}


\newcommand{\sfull}{\textsc{FullSupport}\xspace}
\newcommand{\spart}{\textsc{PartialSupport}\xspace}
\newcommand{\rfull}{\textsc{FullRefute}\xspace}
\newcommand{\rpart}{\textsc{PartialRefute}\xspace}

\newcommand{\scoarse}{\textsc{Support}\xspace}
\newcommand{\rcoarse}{\textsc{Refute}\xspace}

\newcommand{\primary}{\textsc{Primary}\xspace}
\newcommand{\supplemental}{\textsc{Supplemental}\xspace}

\newcommand{\bert}{\textsc{BERT}\xspace}
\newcommand{\scibert}{\textsc{SciBERT}\xspace}
\newcommand{\roberta}{RoBERTa\xspace}
\newcommand{\bioroberta}{BioMedRoBERTa\xspace }
\newcommand{\robertabase}{RoBERTa-base\xspace}
\newcommand{\robertalarge}{RoBERTa-large\xspace}

\newcommand{\softmax}{\textsc{SoftMax}}
\newcommand{\sigmoid}{\textsc{Sigmoid}}

\newcommand{\cls}{\textsc{CLS}\xspace}
\newcommand{\sep}{\textsc{SEP}\xspace}


\newcommand{\multistage}{\textsc{Multistage}\xspace}


\newcommand{\multinli}{\textsc{MultiNLI}\xspace}
\newcommand{\zeroshot}{\textsc{ZeroShot}\xspace}

\newcommand{\nse}{\textsc{NoScoreEv}\xspace}
\newcommand{\se}{\textsc{ScoreEv}\xspace}


\newcommand{\Dcited}{\cD_{\mathrm{cited}}}
\newcommand{\Dseed}{\cD_{\mathrm{seed}}}
\newcommand{\Ddecoy}{\cD_{\mathrm{decoy}}}
\newcommand{\dcited}{\bd_{\mathrm{cited}}}
\newcommand{\dseed}{\bd_{\mathrm{seed}}}
\newcommand{\dciting}{\bd_{\mathrm{citing}}}
\newcommand{\sciting}{\bs_{\mathrm{citing}}}


\newcommand{\fracSupportIncrease}{9\% \xspace}
\newcommand{\fracSupportReduce}{5\% \xspace}
\newcommand{\fracRefuteIncrease}{12\% \xspace}
\newcommand{\fracRefuteReduce}{6\% \xspace}

\newcommand{\bestFeverRetrieval}{0.80 \xspace}
\newcommand{\bestFeverEntailment}{0.90 \xspace}

\newcommand{\feverPaperEntailment}{0.88 \xspace}

\newcommand{\yhat}{\widehat{y}}
\newcommand{\rhat}{\widehat{r}}
\newcommand{\Rhat}{\widehat{R}}
\newcommand{\Shat}{\widehat{S}}
\newcommand{\shat}{\widehat{s}}
\newcommand{\rstar}{r^*}
\newcommand{\Sstar}{S^*}
\newcommand{\Ehat}{\widehat{\cE}}
\newcommand{\Estar}{\cE^*}
\newcommand{\ahat}{\widehat{a}}

\definecolor{NiceGreen}{RGB}{125,204,110}
\definecolor{NiceRed}{RGB}{222,0,0}
\definecolor{c1}{rgb}{0.12, 0.46, 0.70}
\definecolor{c2}{rgb}{1.0, 0.50, 0.05}
\definecolor{c3}{rgb}{0.17,0.63,0.17}

\newcommand{\cmark}{{\color{NiceGreen} \ding{51}}} \newcommand{\xmark}{{\color{NiceRed} \ding{55}}}  

\newcommand{\best}{\textsc{Best}\xspace}

\newcommand{\fonecoarse}{$F1_{\textrm{abstract}}$\xspace}
\newcommand{\foneAbstLab}{$F1_{\textsc{A,L}}$\xspace}
\newcommand{\foneAbstRat}{$F1_{\textrm{A,R}}$\xspace}
\newcommand{\fonefine}{$F1_{\textrm{sentence}}$\xspace}
\newcommand{\foneRatSel}{$F1_{\textrm{R,S}}$\xspace}
\newcommand{\foneRatLab}{$F1_{\textrm{R,L}}$\xspace}

\newcommand{\abstractLabelOnly}{$\textrm{Abstract}_{\textrm{Label-Only}}$\xspace}
\newcommand{\abstractLabelRationale}{$\textrm{Abstract}_{\textrm{Label+Rationale}}$\xspace}
\newcommand{\sentenceSelectionOnly}{$\textrm{Sentence}_{\textrm{Selection-Only}}$\xspace}
\newcommand{\sentenceSelectionLabel}{$\textrm{Sentence}_{\textrm{Selection+Label}}$\xspace}



\newcommand{\website}{\url{https://github.com/allenai/scifact}}

\newcommand{\abst}{\textrm{abst}\xspace}
\newcommand{\rtnl}{\textrm{rtnl}\xspace}
\newcommand{\lbl}{\textrm{lbl}\xspace}


\newcommand{\componentone}{\textsc{AbstractRetrieval}\xspace}
\newcommand{\componenttwo}{\textsc{RationaleSelection}\xspace}
\newcommand{\componentthree}{\textsc{LabelPrediction}\xspace}

\newcommand{\thinmid}{\arrayrulecolor{black!30}\midrule}
\newcommand{\thickmid}{\arrayrulecolor{black}\midrule}

\newcolumntype{L}[1]{>{\raggedright\let\newline\\\arraybackslash\hspace{0pt}}m{#1}}
\newcolumntype{C}[1]{>{\centering\let\newline\\\arraybackslash\hspace{0pt}}m{#1}}
\newcolumntype{R}[1]{>{\raggedleft\let\newline\\\arraybackslash\hspace{0pt}}m{#1}} 



\newcommand{\ba}{{\mathbf{a}}}
\newcommand{\bah}{{\hat{\ba}}}
\newcommand{\ah}{{\hat{a}}}
\newcommand{\Ah}{{\hat{A}}}
\newcommand{\cA}{{\mathcal{A}}}
\newcommand{\at}{{\tilde{a}}}
\newcommand{\bat}{{\tilde{\ba}}}
\newcommand{\At}{{\tilde{A}}}
\newcommand{\bA}{{\mathbf{A}}}

\newcommand{\bb}{{\mathbf{b}}}
\newcommand{\bbt}{{\tilde{\bb}}}
\newcommand{\cB}{{\mathcal{B}}}
\newcommand{\tb}{{\tilde{b}}}
\newcommand{\tB}{{\tilde{B}}}
\newcommand{\hb}{{\hat{b}}}
\newcommand{\hB}{{\hat{B}}}
\newcommand{\bB}{{\mathbf{B}}}

\newcommand{\bc}{{\mathbf{c}}}
\newcommand{\bch}{{\hat{\mathbf{c}}}}
\newcommand{\bC}{{\mathbf{C}}}
\newcommand{\cC}{{\mathcal{C}}}
\newcommand{\ct}{{\tilde{c}}}
\newcommand{\Ct}{{\tilde{C}}}
\newcommand{\ctc}{\ct^{\ast}}

\newcommand{\bd}{{\mathbf{d}}}
\newcommand{\bD}{{\mathbf{D}}}
\newcommand{\cD}{{\mathcal{D}}}
\newcommand{\hd}{{\hat{d}}}  \newcommand{\dt}{{\tilde{d}}}
\newcommand{\bdt}{{\tilde{\bd}}}
\newcommand{\Dt}{{\tilde{D}}}
\newcommand{\dtc}{\dt^{\ast}}

\newcommand{\et}{{\tilde{e}}}
\newcommand{\bfe}{{\mathbf{e}}}
\newcommand{\bE}{{\mathbf{E}}}
\newcommand{\cE}{{\mathcal{E}}}
\newcommand{\cEt}{{\tilde{\cE}}}
\newcommand{\cEb}{{\bar{\cE}}}
\newcommand{\bcE}{{\mathbf{\cE}}}  

\newcommand{\bff}{{{\mathbf{f}}}}
\newcommand{\bF}{{\mathbf{F}}}
\newcommand{\cF}{{\mathcal{F}}}
\newcommand{\ft}{{\tilde{f}}}
\newcommand{\Ft}{{\tilde{F}}}
\newcommand{\Fh}{{\hat{F}}}
\newcommand{\ftc}{\ft^{\ast}}
\newcommand{\bft}{{\tilde{\bff}}}
\newcommand{\bFt}{{\tilde{\bF}}}
\newcommand{\fh}{{\hat{f}}}

\newcommand{\bg}{{\mathbf{g}}}
\newcommand{\gt}{{\tilde{g}}}
\newcommand{\bgt}{{\tilde{\bg}}}
\newcommand{\bG}{{\mathbf{G}}}
\newcommand{\cG}{{\mathcal{G}}}
\newcommand{\Gt}{{\tilde{\bG}}}

\newcommand{\hti}{{\tilde{h}}}
\newcommand{\Hti}{{\tilde{H}}}
\newcommand{\bh}{{\mathbf{h}}}
\newcommand{\bht}{{\tilde{\bh}}}
\newcommand{\Hh}{{\hat{H}}}
\newcommand{\bH}{{\mathbf{H}}}
\newcommand{\bHh}{{\hat{\mathbf{H}}}}

\newcommand{\ih}{{\hat{\imath}}}
\newcommand{\bI}{{\mathbf{I}}}
\newcommand{\cI}{{\mathcal{I}}}

\newcommand{\jh}{{\hat{\jmath}}}
\newcommand{\bJ}{{\mathbf{J}}}
\newcommand{\cJ}{{\mathcal{J}}}
\newcommand{\Jt}{{\tilde{J}}}

\newcommand{\bk}{{\mathbf{k}}}
\newcommand{\bK}{{\mathbf{K}}}
\newcommand{\Kt}{{\tilde{K}}}
\newcommand{\Kh}{{\hat{K}}}
\newcommand{\cK}{{\mathcal{K}}}

\newcommand{\cl}{\ell}
\newcommand{\bL}{{\mathbf{L}}}
\newcommand{\cL}{{\mathcal{L}}}

\newcommand{\mb}{{\mathbf{m}}}
\newcommand{\mh}{{\hat{m}}}
\newcommand{\bM}{{\mathbf{M}}}
\newcommand{\cM}{{\mathcal{M}}}


\newcommand{\cN}{{\mathcal{N}}}
\newcommand{\Nt}{{\tilde{N}}}
\newcommand{\tN}{{\tilde{N}}}  

\newcommand{\bo}{{\mathbf{o}}}
\newcommand{\cO}{{\mathcal{O}}}

\newcommand{\bp}{{\mathbf{p}}}
\newcommand{\bP}{{\mathbf{P}}}
\newcommand{\cP}{{\mathcal{P}}}
\newcommand{\ph}{{\hat{p}}}
\newcommand{\Ph}{{\hat{P}}}

\newcommand{\bq}{{\mathbf{q}}}
\newcommand{\cQ}{{\mathcal{Q}}}
\newcommand{\bQ}{{\mathbf{Q}}}

\newcommand{\br}{{\mathbf{r}}}
\newcommand{\bR}{{\mathbf{R}}}
\newcommand{\cR}{{\mathcal{R}}}
\newcommand{\Rt}{{\tilde{R}}}

\newcommand{\sh}{{\hat{s}}}
\newcommand{\sck}{{\check{s}}}
\newcommand{\shh}{{\Hat{\Hat{s}}}}
\newcommand{\bs}{{\mathbf{s}}}
\newcommand{\bsh}{{\hat{\mathbf{s}}}}
\newcommand{\bsc}{{\check{\mathbf{s}}}}
\newcommand{\bshh}{{\Hat{\Hat{\mathbf{s}}}}}
\newcommand{\bS}{{\mathbf{S}}}
\newcommand{\cS}{{\mathcal{S}}}
\newcommand{\st}{{\tilde{s}}}

\newcommand{\bT}{{\mathbf{T}}}
\newcommand{\bt}{{\mathbf{t}}}
\newcommand{\bu}{{\mathbf{u}}}
\newcommand{\bU}{{\mathbf{U}}}
\newcommand{\bUt}{{\tilde{\bU}}}
\newcommand{\ut}{{\tilde{u}}}
\newcommand{\cU}{{\mathcal{U}}}

\newcommand{\vh}{{\hat{v}}}
\newcommand{\bv}{{\mathbf{v}}}
\newcommand{\bV}{{\mathbf{V}}}
\newcommand{\cV}{{\mathcal{V}}}

\newcommand{\bw}{{\mathbf{w}}}
\newcommand{\bwt}{{\tilde{\bw}}}
\newcommand{\bW}{{\mathbf{W}}}
\newcommand{\cW}{{\mathcal{W}}}

\newcommand{\bx}{{\mathbf{x}}}
\newcommand{\bxt}{{\tilde{\bx}}}
\newcommand{\xt}{{\tilde{x}}}
\newcommand{\Xt}{{\tilde{X}}}
\newcommand{\bX}{{\mathbf{X}}}
\newcommand{\cX}{{\mathcal{X}}}
\newcommand{\bXt}{{\tilde{\bX}}}
\newcommand{\xh}{{\hat{x}}}
\newcommand{\xc}{{\check{x}}}
\newcommand{\xhh}{{\Hat{\Hat{x}}}}
\newcommand{\bxh}{{\hat{\bx}}}
\newcommand{\bxc}{{\check{\bx}}}
\newcommand{\bxhh}{{\Hat{\hat{\bx}}}}

\newcommand{\cY}{{\mathcal{Y}}}
\newcommand{\by}{{\mathbf{y}}}
\newcommand{\byt}{{\tilde{\by}}}
\newcommand{\bY}{{\mathbf{Y}}}
\newcommand{\Yt}{{\tilde{Y}}}
\newcommand{\yt}{{\tilde{y}}}
\newcommand{\yh}{{\hat{y}}}

\newcommand{\zt}{{\tilde{z}}}
\newcommand{\zh}{{\hat{z}}}
\newcommand{\bz}{{\mathbf{z}}}
\newcommand{\bZ}{{\mathbf{Z}}}
\newcommand{\cZ}{{\mathcal{Z}}}



\newcommand{\al}{\alpha}
\newcommand{\bal}{{\boldsymbol{\al}}}
\newcommand{\balh}{{\hat{\boldsymbol{\al}}}}
\newcommand{\alh}{{\hat{\al}}}

\newcommand{\btt}{{\tilde{\bt}}}
\newcommand{\btht}{{\hat{\bt}}}

\newcommand{\g}{\gamma}
\newcommand{\G}{\Gamma}
\newcommand{\bGa}{{\boldsymbol{\Gamma}}}
\newcommand{\gh}{{\hat{\g}}}

\newcommand{\de}{\delta}
\newcommand{\De}{\Delta}
\newcommand{\Deh}{{\hat{\Delta}}}
\newcommand{\bde}{{\boldsymbol{\de}}}
\newcommand{\bDe}{{\boldsymbol{\De}}}

\newcommand{\e}{\epsilon}
\newcommand{\eps}{\varepsilon}

\newcommand{\etah}{{\hat{\eta}}}
\newcommand{\bpi}{{\boldsymbol{\pi}}}

\newcommand{\pht}{{\tilde{\phi}}}
\newcommand{\Pht}{{\tilde{\Phi}}}

\newcommand{\pst}{{\tilde{\psi}}}
\newcommand{\Pst}{{\tilde{\Psi}}}

\newcommand{\s}{\sigma}
\newcommand{\sih}{\hat{\sigma}}

\newcommand{\z}{\zeta}
\newcommand{\ztt}{{\tilde{\z}}}
\newcommand{\ztb}{{\bar{\z}}}

\newcommand{\thh}{{\hat{\theta}}}
\newcommand{\Thh}{{\hat{\Theta}}}
\newcommand{\Th}{\Theta}
\newcommand{\bth}{{\boldsymbol{\theta}}}
\newcommand{\bTh}{{\boldsymbol{\Theta}}}
\newcommand{\bThh}{{\hat{\bTh}}}
\newcommand{\Tht}{{\tilde{\Theta}}}

\newcommand{\la}{\lambda}
\newcommand{\La}{\Lambda}
\newcommand{\lam}{\lambda}  \newcommand{\Lam}{\Lambda}  \newcommand{\bLa}{{\boldsymbol{\La}}}
\newcommand{\lah}{{\hat{\lam}}}

\newcommand{\bmu}{{\boldsymbol{\mu}}}

\newcommand{\bXi}{{\boldsymbol{\Xi}}}

\newcommand{\bPi}{{\boldsymbol{\Pi}}}

\newcommand{\rht}{{\tilde{\rho}}}
\newcommand{\rhc}{{\check{\rho}}}

\newcommand{\bSi}{{\boldsymbol{\Sigma}}}

\newcommand{\ups}{\upsilon}
\newcommand{\Ups}{\Upsilon}
\newcommand{\bUp}{{\boldsymbol{\Ups}}}

\newcommand{\bPs}{{\boldsymbol{\Psi}}}

\newcommand{\w}{\omega}
\newcommand{\wh}{{\hat{\omega}}}
\newcommand{\W}{\Omega}
 
\title{Fact or Fiction: Verifying Scientific Claims}

\author{David Wadden$^\mathbf{\dagger}$\thanks{\xspace \xspace Work performed during internship with the Allen Institute for Artificial Intelligence} \quad
Shanchuan Lin$^\mathbf{\dagger}$ \quad
Kyle Lo$^\mathbf{\ddagger}$ \quad
Lucy Lu Wang$^\mathbf{\ddagger}$ \\
{\bf Madeleine van Zuylen}$^\mathbf{\ddagger}$ \quad
{\bf Arman Cohan}$^\mathbf{\ddagger}$ \quad
{\bf Hannaneh Hajishirzi}$^\mathbf{\dagger\ddagger}$\\
  $^\mathbf{\dagger}$
  University of Washington, Seattle, WA, USA \\
  $^\mathbf{\ddagger}$
  Allen Institute for Artificial Intelligence, Seattle, WA, USA \\
  {\tt\small \{dwadden,linsh,hannaneh\}@cs.washington.edu} \\
  {\tt\small \{kylel,lucyw,madeleinev,armanc\}@allenai.org}
}

\date{}

\begin{document}
\maketitle


\begin{abstract}






We introduce scientific claim verification, a new task to select abstracts from the research literature containing evidence that supports or refutes a given scientific claim, and to identify rationales justifying each decision. To study this task, we construct \ours, a dataset of 1.4K expert-written scientific claims paired with evidence-containing abstracts annotated with labels and rationales. We develop baseline models for \ours, and demonstrate that these models benefit from combined training on a large dataset of claims about Wikipedia articles, together with the new \ours data. We show that our claim verification system is able to identify plausible evidence for 23 / 36 claims relevant to COVID-19 on the CORD-19 corpus. Our results and experiments strongly suggest that our new task and data will support significant future research efforts.\footnote{Data, code, and a web demo will released publicly.}


\end{abstract} \section{Introduction} \label{sec:introduction}




Due to rapid growth in the scientific literature, it is difficult for researchers -- and the general public even more so -- to stay up to date on the latest findings. This challenge is especially acute during public health crises like the current COVID-19 pandemic, due to the extremely fast rate at which new findings are reported and the risks associated with making decisions based on outdated or incomplete information. As a result, there is a need for automated tools to assist researchers and the public in evaluating the veracity of scientific claims. 


\emph{Fact-checking} -- a task in which the veracity of an input \emph{claim} is verified against a corpus of documents that \emph{support} or \emph{refute} the claim -- has been studied to combat the proliferation of misinformation in political news, social media, and on the web~\cite{Thorne2018FEVERAL,Hanselowski2019ARA}. However, verifying scientific claims poses new challenges to both dataset construction and effective modeling. While political claims are readily available on fact-checking websites and can be verified by crowd workers, annotators with extensive domain knowledge are required to generate and verify scientific claims. In addition, NLP systems for scientific claim verification must possess expert background knowledge, as well as language understanding and reasoning capabilities not required to verify factoid claims. 

\begin{figure}[t]
  \centering

  \includegraphics[width=\columnwidth]{figures/teaser-fig.pdf}
  \caption{A \ours claim refuted by evidence. To refute this claim, the system must recognize that (1) ``CSF'' is an acronym for ``cerebral spinal fluid'', found in the brain, (2) ``Citalopram'' is a type of antidepressant, but ``placebo'' is not, and (3) ``Slowing by 37\%'' indicates a reversal in effect relative to the claim.}
  \label{fig:teaser}
\end{figure} \definecolor{c1}{RGB}{84,130,53}
\newcommand{\supportsCovid}[1]{{\color{c1} {#1}}}

\definecolor{c2}{RGB}{218,0,0}
\newcommand{\refutesCovid}[1]{{\color{c2} {#1}}}

\definecolor{c3}{RGB}{127,127,127}
\newcommand{\wrongContext}[1]{{\color{c3} \textbf{#1}}}



\begin{table*}[t!]
  \setlength{\tabcolsep}{.25em}
  \footnotesize
  \centering

  \begin{tabularx}{\textwidth}{ X }
    \toprule
    \textbf{Claim 1}: Lopinavir / ritonavir have exhibited favorable clinical responses when used as a treatment for coronavirus. \\
    \midrule
    \textbf{Supports}: \dots \supportsCovid{Interestingly, after lopinavir/ritonavir (Kaletra, AbbVie) was administered, $\beta$-coronavirus viral loads significantly decreased and no or little coronavirus titers were observed.} \\
    \midrule
    \textbf{Refutes}: \refutesCovid{The focused drug repurposing of known approved drugs (such as lopinavir/ritonavir) has been reported failed for curing SARS-CoV-2 infected patients.}. It is urgent to generate new chemical entities against this virus \dots \\
    \midrule
    \textbf{Claim 2}: The coronavirus cannot thrive in warmer climates. \\
    \midrule
    \textbf{Supports}: \supportsCovid{...most outbreaks display a pattern of clustering in relatively cool and dry areas...This is because the environment can mediate human-to-human transmission of SARS-CoV-2, and unsuitable climates can cause the virus to destabilize quickly...} \\
    \midrule
    \textbf{Refutes}: \refutesCovid{...significant cases in the coming months are likely to occur in more humid (warmer) climates, irrespective of the climate-dependence of transmission and that summer temperatures will not substrantially limit pandemic growth}. \\
    \bottomrule
  \end{tabularx}

  \caption{Evidence identified by our system as supporting and refuting two claims concerning COVID-19.}
  \label{tbl:covid_example}


\end{table*}
 In this paper, we introduce the task of scientific claim verification (Figure~\ref{fig:teaser}) to evaluate the veracity of scientific claims against a research corpus. Table \ref{tbl:covid_example} presents some examples. To facilitate research on this task, we construct \ours, an expert-annotated dataset of \numclaims scientific claims accompanied by scientific abstracts that support or refute each claim, and annotated with rationales \cite{Lei2016RationalizingNP} justifying each support / refute decision. To create the dataset, we develop a novel annotation protocol in which annotators re-formulate naturally occurring claims in the scientific literature -- \emph{citation sentences} -- into atomic scientific claims. Using citation sentences as a source of claims both speeds the claim generation process and guarantees that the topics discussed in \ours are representative of the research literature. In addition, citation links indicate the exact documents likely to contain evidence necessary to verify a given claim.

We establish performance baselines on \ours following \cite{DeYoung2019ERASERAB}, which achieves strong performance on \fever, a Wikipedia fact checking dataset~\cite{Thorne2018FEVERAL}. Our baseline is a pipeline system which retrieves abstracts related to an input claim, uses a \bert-based \cite{Devlin2019BERTPO} sentence selector to identify rationale sentences, and labels each abstract as \supports, \refutes, or \notenough with respect to the claim. We show that our baseline can benefit from training on fact-checking data from domains including Wikipedia articles and politics.  



We showcase the ability of our model to verify expert-written claims concerning the novel coronavirus COVID-19 against the newly-released CORD-19 corpus \cite{Wang2020CORD19TC}. Expert annotators judge retrieved evidence to be plausible for 23 of 36 claims.\footnote{We emphasize that our model is a \emph{research prototype} and should not be used to make any medical decisions whatsoever.} Our results and analyses demonstrate the importance of the new task and dataset to support significant future research in this domain.
%
 \section{Background and task definition}

As illustrated in Figure \ref{fig:teaser}, scientific claim verification  is the task of identifying evidence from the research literature that \supports or \refutes a given scientific claim. Table \ref{tbl:covid_example} shows the results of our system  applied to claims about the novel coronavirus COVID-19. For each claim, the system identifies relevant scientific abstracts, and labels the relation of each abstract to the claim as either \supports or \refutes. Verifying scientific claims is challenging and requires domain-specific background knowledge -- for instance, in order to identify the evidence supporting Claim 1 in  Table~\ref{tbl:covid_example}, the system must determine that a reduction in coronavirus viral load indicates a favorable clinical response, even though this fact is never mentioned. 

\vspace{.2cm}
\noindent{\bf Scientific claims} \label{sec:claims} In \ours, a scientific claim is an \emph{atomic verifiable statement} expressing a finding about one aspect of a scientific entity or process, which can be verified from a single source.\footnote{Requiring annotators to search multiple sources increases cognitive burden and decreases annotation quality.}
For instance,  ``\emph{The $R_0$ of the novel coronavirus is 2.5}'' is valid, but opinion-based statements like ``\emph{The government should require people to stand six feet apart to stop coronavirus}'' are not. Compound claims like ``\emph{Aerosolized coronavirus droplets can travel at least 6 feet and can remain in the air for 3 hours}'' should be split into two atomic claims. 

Claims in \ours  are \emph{natural} -- they are derived from citation sentences, or \emph{citances} \cite{Nakov2004CitancesC}, that occur naturally in scientific articles. This is similar to political fact-checking datasets such as \snopes \cite{Hanselowski2019ARA}, which use political fact-checking websites as a source of natural claims. On the other hand, claims in the popular \fever dataset \cite{Thorne2018FEVERAL} are \emph{synthetic}, since they are created by annotators by mutating sentences from the Wikipedia articles that will serve as evidence.

\vspace{.2cm}
\noindent{\bf Supporting and refuting evidence} \label{sec:evidence}
In most fact-checking work, claims are assigned a global supported or refuted label based on the entirety of the available evidence. For example in \fever, the claim ``\emph{Barack Obama was the $44^{th}$ President of the United States}'' can be verified as supported given Wikipedia as a corpus of evidence.

While \ours claims are indeed verifiable assertions about scientific findings, accurately assigning a global truth label to a scientific claim (given a fixed scientific corpus) requires a systematic review by a team of experts. In this work we focus on the simpler task of assigning \supports or \refutes relations to individual \emph{claim-abstract pairs}.

Each \supports or \refutes relation between claim and abstract must be justified by at least one \emph{rationale}. A rationale is a minimal collection of sentences which, taken together as premises in the context of the abstract, can reasonably be judged by a domain expert as implying the claim. Rationales facilitate the development of \emph{interpretable models} which not only have the ability to make label predictions, but can also identify the exact sentences that are necessary for their decisions.






%
 \section{The \ours dataset} \label{sec:dataset}


\label{sec:dataset_overview}

The \ours dataset consists of \numclaims scientific claims verified against a corpus of \numpapers abstracts.  Abstracts that support or refute each claim are annotated with rationales.  We describe our corpus creation and annotation process. Additional details on the corpus and data collection process can be found in Appendix \ref{sec:data_collection}, along with additional statistics on the dataset.



\subsection{Data source and corpus construction}
\label{sec:data_source}

To construct \ours, we use S2ORC \cite{Lo2019GORCAL}, a publicly-available corpus of millions of scientific articles. To ensure that documents in our dataset are of high quality, we randomly sample articles from a manually curated collection of well-regarded journals spanning domains from basic science (e.g., \emph{Cell}, \emph{Nature}) to clinical medicine (e.g., \emph{JAMA}, \emph{BMJ}). We restrict to articles with at least 10 citations and with full text freely available. The resulting collection is referred to as our \textit{seed} set. We use the S2ORC citation graph to sample \emph{source citances} from \emph{citing articles} which cite these seed articles. If a citance cites other articles not in the seed set, we refer to these as \emph{co-cited} articles and add them to the corpus, as depicted in Figure \ref{fig:corpus}. The majority of citances used for \ours cite only the seed article (no co-cited articles), as we found in initial annotation experiments that these citances tended to yield specific, easy-to-verify claims.






To expand the corpus, we identify five papers cited in the same paper as each source citance but in a different paragraph, and add these to the corpus as \emph{distractor abstracts}. These abstracts often discuss similar topics to the evidence documents, increasing the difficulty of abstract retrieval and making our metrics more accurately reflect the system's performance on a large research corpus.

\begin{figure}[t]
  \centering

  \includegraphics[width=\columnwidth, height=6.5cm]{figures/corpus-fig.pdf}
  \caption{Corpus construction. Citing abstracts are identified for each seed document. A claim is written based on the source citance in the citing abstract.}
  \label{fig:corpus}

\end{figure} 


\subsection{Claim writing} \label{sec:claim_writing}
\noindent {\bf Annotation}
Annotators are shown a source citance in the context of an article, and are asked to write up to three claims based on the content of the citance. This results in \emph{natural} claims because the annotator does not see the cited article's abstract -- the \emph{cited abstract} -- at the time of claim writing. They are asked to skip citances that do not make statements about specific scientific findings. All annotators for both claim writing and verification were advanced undergraduate or graduate students with training in the life sciences.


\vspace{.2cm} \noindent{\bf Claim negation} Unless the authors of the source citance were mistaken, cited articles should provide supporting evidence for the claims made in a citance. To obtain examples where an abstract \refutes a claim, a scientific NLP expert wrote negations of existing claims, taking precautions not to bias the negations by relying on obvious keywords like ``not'' \cite{Schuster2019TowardsDF}.




\subsection{Claim verification}
\label{sec:claim_verification}

\noindent {\bf Annotation} For each claim, all of the claim's cited abstracts are annotated for evidence. Annotators are shown a single claim / cited abstract pair, and asked to label the pair as \supports, \refutes, or \notenough. Although our task definition allows for a single claim to be both supported and refuted (by different abstracts) -- an occurrence we observe on real-world COVID-19 claims (\S\ref{sec:covid}) -- this never occurs in our dataset. Each claim has a single label. Counts for each label are shown in Table \ref{tbl:label_counts}. Overall, the annotators found evidence in 63\% of cited abstracts.  If the annotator assigns a \supports or \refutes label, they must also identify all rationales as defined in \S\ref{sec:evidence}. Table \ref{tbl:evidence_counts} provides statistics on the number of sentences per rationale, the number of rationales per claim / abstract pair, and the number of evidence abstracts per claim.  No abstract has more than 3 rationales for a given claim, and all rationales consist of at most three sentences. Rationales in \ours are mutually exclusive. 28 rationales contain non-contiguous sentences.

\ours claims are verified against abstracts rather than full articles since (1) abstracts can be annotated more scalably, (2) evidence is found in the abstract in more than 60\% of cases, and (3) previous attempts at full-document annotation suffered from low annotator agreement (\S\ref{sec:related_scinlp}).

\begin{table}[t]
  \setlength{\tabcolsep}{.25em}
  \footnotesize
  \centering


  \begin{subtable}{\linewidth}
    \begin{tabular}{L{0.12\linewidth} C{0.19\linewidth} C{0.19\linewidth} C{0.19\linewidth} C{0.15\linewidth}}
      \toprule
      Fold &  \supports &  \notenough &  \refutes &   All \\
      \midrule
      Train &       332 &              304 &      173 &   809 \\
      Dev   &       124 &              112 &       64 &   300 \\
      Test  &       100 &              100 &      100 &   300 \\
      \midrule
      All   &       556 &              516 &      337 &  1409 \\
      \bottomrule
    \end{tabular}
    \subcaption{Distribution of claim labels in \ours.}
    \label{tbl:label_counts}
  \end{subtable}

  \vspace{0.25em}

  \begin{subtable}{\linewidth}
    \begin{tabular}{L{0.48\linewidth} C{0.1\linewidth} C{0.1\linewidth} C{0.1\linewidth} C{0.1\linewidth}}
      \toprule
      {} &    0 &     1 &    2 &   3+ \\
      \midrule
      Cited abstracts per claim     &    - &  1278 &   86 &   45 \\
      Evidence abstracts per claim &  516 &   830 &   37 &   26 \\
      Rationales per abstract         &    - &   552 &  290 &  153 \\
      Sentences per rationale       &    - &  1542 &   92 &   11 \\
      \bottomrule
    \end{tabular}
    \subcaption{Evidence counts at various levels of granularity. For example, Column 2 of the row ``Rationales / abstract'' indicates that 290 claim / abstract pairs are supported by 2 distinct rationales.}
    \label{tbl:evidence_counts}
  \end{subtable}

  \caption{Statistics on claim labels, and the number of evidence abstracts and rationales per claim.}
  \label{tbl:evidence_stats}


\end{table} 
\vspace{.2cm} \noindent{\bf Quality}
We assign \nagreement claim-abstract pairs for independent re-annotation. The label agreement is \labelkappa Cohen's $\kappa$, comparable with the 0.68 Fleiss' $\kappa$ reported in \cite{Thorne2018FEVERAL}, and 0.70 Cohen's $\kappa$ reported in \cite{Hanselowski2019ARA}. To measure rationale agreement, we treat each sentence as either classified as ``part of a rationale'' or ``not part of a rationale'' and compute sentence-level agreement on abstracts where annotators agreed on the \entailment label. The resulting Cohen's $\kappa$ is \rationalekappaagree.
 \section{The \ours task} \label{sec:task_definition}
\vspace{.2cm} \noindent {\bf Task Formulation} The inputs to our task are a scientific claim $c$ and a corpus of abstracts $\cA$.
All abstracts $a \in \cA$ are labeled as $y(c, a)\in\{$\supports, \refutes, \notenough$\}$ with respect to a claim $c$.  The abstracts that either \support or \refute $c$ are referred to as \emph{evidence abstracts} for $c$, denoted as $\cE(c)$. Each evidence abstract $a \in \cE(c)$ is annotated with rationales. A single rationale $R_i$ is a collection of sentences $\{r_1(c, a), \dots, r_{m}(c, a)\}$, where $m$ is the number of sentences in rationale $R_i$.
We denote the set of all rationales as $\cR(c, a) = \{R_1(c, a), \dots, R_n(c, a)\}$.

Given a claim $c$ and a corpus $\cA$, the system must predict a set of evidence abstracts $\Ehat(c)$. For each abstract $a \in \Ehat(c)$, it must predict a label $\yhat(c, a)$, and a collection of \emph{rationale sentences} $\Shat(c, a) = \{\shat_1(c, a), \dots, \shat_{\ell}(c, a)\}$. Note that although the gold annotations may contain multiple separate rationales, to simplify the prediction task we only require the model to predict a single collection of \emph{rationale sentences}; these sentences may encompass multiple gold rationales.


\vspace{.2cm} \noindent {\bf Task Evaluation} \label{sec:evaluation}
We evaluate the task at two levels of granularity. For \emph{abstract-level} evaluation, we assess the model's ability to identify the abstracts in the corpus that support or refute the claim. For \emph{sentence-level} evaluation, we evaluate the model's performance at identifying the sentences sufficient to justify the abstract-level predictions. We conduct evaluations in both the ``Open'' \fever-style \cite{Thorne2018FEVERAL} setting where the evidence abstracts must be retrieved, and the ``Oracle abstract'' \eraser-style \cite{DeYoung2019ERASERAB} setting where the gold evidence abstracts $\cE(c)$ are provided.


\vspace{.2cm} \noindent {\bf Abstract-level evaluation} is inspired by the \fever score. Given a claim $c$, a predicted evidence abstract $a \in \Ehat(c)$ is \emph{correctly labeled} if (1) $a$ is a gold evidence abstract for $c$, and (2) The predicted label is correct: $\yhat(c, a) = y(c, a)$. It is \emph{correctly rationalized} if, in addition, the predicted rationale sentences contain a gold rationale, i.e., there exists some gold rationale $R_i(c, a) \subseteq \Shat(c, a)$.

Like \fever, which limits the maximum number of predicted rationale sentences to five, \ours limits to three predicted rationale sentences. Overall performance is measured by the micro-F1 of the precision and recall over the correctly-labeled and correctly-rationalized evidence abstracts. We refer to these evaluations as \abstractLabelOnly and \abstractLabelRationale, respectively.

\vspace{.2cm} \noindent {\bf Sentence-level evaluation} measures performance in identifying individual rationale sentences. Unlike the abstract-level metrics, this evaluation penalizes the prediction of extra rationale sentences. A predicted rationale sentence $\shat(c, a)$ is correctly \emph{selected} if (1) It is a member of some gold rationale $R_i(c, a)$ and (2) All other sentences from the same gold rationale $R_i(c, a)$ are among the predicted $\Shat(c, a)$. It is correctly \emph{labeled} if, in addition, the abstract $a$ is correctly labeled: $\yhat(c, a) = y(c, a)$. Overall performance is measured by the micro-F1 of the precision and recall of correctly-selected and correctly-labeled rationale sentences, denoted \sentenceSelectionOnly and \sentenceSelectionLabel, respectively. For sentence-level evaluation, we do not limit the number of predicted rationale sentences, since the evaluation penalizes models that over-predict.
 \section{\sysname: Baseline model} \label{sec:baselines}
We develop a baseline (referred to as \sysname) that takes a claim $c$ and corpus $\cA$ as input, identifies evidence abstracts $\Ehat(c)$, and predicts a label $\yhat(c,a)$ and rationale sentences $\Shat(c, a)$ for each $a \in \Ehat(c)$.  Following the ``\bert-to-\bert''  model presented in \cite{DeYoung2019ERASERAB,Soleimani2019BERTFE}, \sysname\ is a pipeline of three components:

\begin{enumerate}[leftmargin=*]
    \item \componentone retrieves $k$ abstracts with highest TF-IDF similarity to the claim.
    \item \componenttwo identifies rationale sentences $\Shat(c, a)$ for each abstract.
    \item \componentthree makes the final label prediction $\yhat(c, a)$.
\end{enumerate}



\vspace{.2cm} \noindent {\bf Rationale selection} \label{sec:rationale_selection}
Given a claim $c$ and abstract $a$, we train a model to predict $z_i \triangleq \mathds{1}[a_i \textrm{ is a rationale sentence}]$ for each sentence $a_i$ in $a$. For each sentence, we encode the concatenated sequence $w_i = [a_i, \sep, c]$ using a \bert-style language model and predict a score $\tilde{z}_i = \sigma[f(\cls(w_i))]$, where $\sigma$ is the sigmoid function, $f$ is a linear layer and $\cls(w_i)$ is the \cls token from the encoding of $w_i$. We minimize cross-entropy loss between $z_i$ and $\tilde{z}_i$ during training. We train the model on pairs of claims and their cited abstracts. For each claim, we use cited abstracts labeled \notenough, as well as non-rationale sentences from abstracts labeled \supports and \refutes as negative examples. To make predictions, we select all sentences $a_i$ with $\tilde{z_i} > t$ as rationale sentences, where $t \in [0, 1]$ is tuned on the dev set (Appendix \ref{sec:parameters}).





\vspace{.2cm} \noindent {\bf Label prediction} \label{sec:entailment}
Sentences identified by the rationale selector are passed to a separate \bert-based model to make the final
labeling decision. Given a claim $c$ and abstract $a$, we concatenate the claim and the predicted rationale sentences $u = [\shat_1(c, a), \dots \shat_{\ell}(c, a), \sep, c]$\footnote{We truncate the rationale input if it exceeds the \bert token limit. $c$ is never truncated.}, and predict $\tilde{y}(c, a) = \phi[f(\cls(u))]$, where $\phi$ is the softmax function, and $f$ is a linear layer with three outputs representing the \{\supports, \refutes, \notenough\} labels. We minimize the cross-entropy loss between $\tilde{y}(c, a)$ and the true label $y(c, a)$.

We train the model on pairs of claims and their cited abstracts using gold rationales as input. For cited abstracts labeled \notenough, we choose the $k$ sentences from the cited abstract with highest TF-IDF similarity to the claim as input rationales. For prediction, we use the predicted rationale sentences $\Shat(c, a)$ as input and predict $\hat{y}(c, a) = \textrm{argmax }\tilde{y}(c, a)$. \notenough is predicted for abstracts with no rationale sentences.

We experimented with a label prediction model which encodes entire abstracts via the Longformer \cite{Beltagy2020LongformerTL}, and makes predictions using the document-level \cls token. Performance was not competitive with our pipeline setup, likely because the label predictor struggles to identify relevant information when given full abstracts. \section{Experiments}

In our experiments, we (1) analyze the performance of each individual component of \sysname, (2) evaluate full task performance in both the ``Oracle abstract'' and ``Open'' settings, (3) present promising results verifying claims about COVID-19 using \sysname, and (4) discuss some modeling challenges presented by the dataset.









\subsection{Pipeline components} \label{sec:pipeline_components}
We examine the effects of different training datasets, sentence encoders, and model inputs on the performance of the \componenttwo and \componentthree modules. See Appendix \ref{sec:model_details} for implementation details.


\begin{table}[t]
  \footnotesize
  \setlength{\tabcolsep}{0.5em}
  \centering

  \begin{tabularx}{\columnwidth}{l *{3}{c}  *{3}{c}}
    \toprule
 & \multicolumn{3}{c}{\textsc{Rational-Select.}} & \multicolumn{3}{c}{\textsc{Label-Pred.}} \\


    \midrule
    \textbf{Training data} & P & R & F1 & \multicolumn{3}{c}{\textsc{Acc.}} \\

    \arrayrulecolor{black!30}\midrule


    \fever         &      41.5 &   57.9 &  48.4 &  \multicolumn{3}{c}{67.6} \\
    \snopes        &      42.5 &   62.3 &  50.5 &  \multicolumn{3}{c}{71.3} \\
    \ours          &      73.7 &   70.5 &  \textbf{72.1} &  \multicolumn{3}{c}{75.7} \\
    \fever + \ours &      72.4 &   67.2 &  69.7 &  \multicolumn{3}{c}{\textbf{81.9}} \\

    \arrayrulecolor{black}
    \midrule

    \textbf{Sentence encoder} & P & R & F1 & \multicolumn{3}{c}{\textsc{Acc.}} \\
    \arrayrulecolor{black!30}\midrule


    \scibert      &      74.5 &   74.3 &  \textbf{74.4} &  \multicolumn{3}{c}{69.2} \\
    \bioroberta   &      75.3 &   69.9 &  72.5 &  \multicolumn{3}{c}{71.7} \\
    \robertabase  &      76.1 &   66.1 &  70.8 &  \multicolumn{3}{c}{62.9} \\
    \robertalarge &      73.7 &   70.5 &  72.1 &  \multicolumn{3}{c}{\textbf{75.7}} \\

    \arrayrulecolor{black}
    \midrule

    \textbf{Model inputs} & P & R & F1 & \multicolumn{3}{c}{\textsc{Acc.}} \\
    \arrayrulecolor{black!30}\midrule
Claim-only         &      - &   - &  - &  \multicolumn{3}{c}{44.5} \\
    Abstract-only         &   60.1 &   60.9 &  60.5 &  \multicolumn{3}{c}{53.3} \\

    \arrayrulecolor{black}
    \bottomrule
  \end{tabularx}

  \caption{Comparison of different training datasets, encoders, and model inputs for \componenttwo and \componentthree, evaluated on the \ours dev set. A claim-only model cannot select rationales.}
  \label{tbl:components}

\end{table} \vspace{.2cm} \noindent{\bf Training Data}
We train on (1) \fever, (2) \snopes, (3) \ours, and (4) \fever pretraining followed by \ours fine-tuning. \robertalarge \cite{Liu2019RoBERTaAR} is used as the sentence encoder.

\vspace{.2cm} \noindent{\bf Sentence encoder}
We fine-tune \scibert \cite{Beltagy2019SciBERTPC}, \bioroberta \cite{Gururangan2020DONTST}, \robertabase, and \robertalarge. \ours is used as training data.

\vspace{.2cm} \noindent{\bf Model Inputs}
We examine the performance of ``claim-only'' and ``abstract-only'' models trained on \ours, using \robertalarge as the sentence encoder. The claim-only model makes label predictions based on the claim text alone, without access to evidence abstracts. The abstract-only model selects rationale sentences and makes label predictions without access to the claim.

\begin{table*}[t]
  \footnotesize
  \setlength{\tabcolsep}{0.5em}
  \centering

  \begin{tabularx}{\linewidth}{*{3}{l} *{6}{c} *{1}{c} *{6}{c}}

    \toprule

    & & & \multicolumn{6}{c}{\textbf{Sentence-level}} & & \multicolumn{6}{c}{\textbf{Abstract-level}} \\

     & & & \multicolumn{3}{c}{\textbf{Selection-Only}} & \multicolumn{3}{c}{\textbf{Selection+Label}} & & \multicolumn{3}{c}{\textbf{Label-Only}} & \multicolumn{3}{c}{\textbf{Label+Rationale}}  \\

    \cmidrule(lr){4-6} \cmidrule(lr){7-9} \cmidrule(lr){11-13} \cmidrule(lr){14-16}

    Retrieval & Model & Row & P & R & F1 & P & R & F1 & & P & R & F1 & P & R & F1 \\

    \midrule

    \multirow{4}{*}{\shortstack[c]{\textbf{Oracle}\\ \textbf{abstract}}} &


    Oracle rationale & 1 & 100.0 & 80.5 & 89.2 & 89.6 & 72.2 & 79.9 &  & 90.1 & 77.5 & 83.3 & 90.1 & 77.5 & 83.3 \\

    \arrayrulecolor{black!30}\cmidrule(lr){2-16}

 & Zero-shot & 2 & 42.5 & 45.1 & 43.8 & 36.1 & 38.4 & 37.2 &  & 86.9 & 53.6 & 66.3 & 67.9 & 41.9 & 51.8 \\

 & \sysname & 3 & 76.1 & 63.8 & 69.4 & 66.5 & 55.7 & \textbf{60.6}          &  & 87.3 & 65.3 & 74.7 & 84.9 & 63.5 & \textbf{72.7}          \\

\thickmid

\multirow{4}{*}{\textbf{Open}} &

Oracle rationale & 4 & 100.0 & 56.5 & 72.2 & 87.6 & 49.5 & 63.2 &  & 88.9 & 54.1 & 67.2 & 88.9 & 54.1 & 67.2 \\

\arrayrulecolor{black!30}\cmidrule(lr){2-16}
 & Zero-shot & 5 & 28.7 & 37.6 & 32.5 & 23.7 & 31.1 & 26.9 &  & 56.0 & 42.3 & 48.2 & 42.3 & 32.0 & 36.4 \\

 & \sysname & 6 & 45.0 & 47.3 & 46.1 & 38.6 & 40.5 & \textbf{39.5}          &  & 47.5 & 47.3 & 47.4 & 46.6 & 46.4 & \textbf{46.5}          \\


\arrayrulecolor{black}\bottomrule

  \end{tabularx}

  \caption{Test set performance on \ours, according to the metrics from \S\ref{sec:evaluation}. For the \emph{Oracle abstract} rows, the system is provided with gold evidence abstracts. ``Oracle rationale'' rows indicate that the gold rationales are provided as input. ``Zero-shot'' indicates zero-shot performance of a verification system trained on \fever. See Appendix \ref{sec:uncertainty} for standard deviations on all estimates, and dev set performance. }
    \label{tbl:main_results}

\end{table*}
 
The \componenttwo module is evaluated on its ability to select rationale sentences given gold abstracts\footnote{Our \fever-trained \componenttwo module achieves 79.9 sentence-level F1 on the \fever test set, virtually identical to 79.6 reported in \cite{DeYoung2019ERASERAB}.}. The \componentthree module is evaluated on its 3-way label classification accuracy given gold rationales from cited abstracts. Cited abstracts labeled \notenough are included in the evaluation. These abstracts have no gold rationale sentences; as in \S\ref{sec:baselines}, we provide the $k$ most similar sentences from the abstract as input.

\vspace{.2cm} \noindent{\bf Results} The results are shown in Table \ref{tbl:components}. For \componentthree, the best performance is achieved by training first on the large \fever dataset and then fine-tuning on the smaller in-domain \ours training set. To understand the benefits of \fever pretraining, we examined the claim / evidence pairs where the \fever + \ours - trained model made correct predictions but the \ours - trained model did not. In  36 / 44 of these cases, the \ours - trained model predicts \notenough. Thus pretraining on \fever appears to improve the model's ability to recognize textual entailment relationships between evidence and claim -- particularly relationships indicated by non-domain-specific cues like ``is associated with'' or ``has an important role in''.

For \componenttwo, training on \ours alone produces the best results. We examined the rationales that the \ours - trained model identified but the \fever - trained model missed, and found that they generally contain science-specific vocabulary. Thus, training on additional out-of-domain data provides little benefit.

\robertalarge exhibits the strongest performance on label prediction, while \scibert has a slight edge on rationale selection. The ``claim-only'' and `'abstract-only'' models perform substantially worse than the model with access to the full input, providing reassurance that the model is not making its predictions based on statistical artifacts.



\subsection{Full task} \label{sec:main_results}
\noindent {\bf Experimental setup}
Based on the results from \S\ref{sec:pipeline_components}, we use the \componenttwo module trained on \ours only, and the \componentthree module trained on \fever + \ours for our final end-to-end system \sysname. Although \scibert performs slightly better on rationale selection, using \robertalarge for both \componenttwo and \componentthree gave the best full-pipeline performance on the dev set, so we use \robertalarge for both components. For the \componentone module, the best dev set full-pipeline performance was achieved by retrieving the top $k=3$ documents. 

\vspace{.2cm}\noindent{\bf Model comparisons}
We report performance of three model variants. For the ``Oracle rationale'' setting, the \componenttwo module is replaced by an oracle which outputs gold rationales for correctly retrieved documents, and no rationales for incorrect retrievals. The ``Zero-shot'' setting reports the zero-shot generalization performance of a model trained on \fever (the results on \snopes were slightly worse). \sysname reports the performance of our best system.


\vspace{.2cm}\noindent{\bf Results}
The results are shown in Table \ref{tbl:main_results}. In the oracle abstract setting, the abstract-level F1 scores are roughly comparable to label classification accuracies, and the \abstractLabelRationale score in Row 3 implies an end-to-end classification accuracy of roughly 70\%, given gold abstracts.

Access to in-domain data during training clearly improves performance. Despite the small size of \ours, training on these data leads to relative improvements of 47\% on open \sentenceSelectionLabel, and 28\% on open \abstractLabelRationale (Row 6) over \fever alone.
The three pipeline components make similar contributions to the overall model error. Replacing \componenttwo with an oracle leads to a roughly 20-point rise in \sentenceSelectionLabel F1 (Row 6 vs. Row 4). Replacing \componentone with an oracle as well leads to a gain of roughly 20 more points (Row 4 vs. Row 1).

Nearly all correctly-labeled abstracts are supported by at least one rationale. There is only a two-point difference in F1 between \abstractLabelOnly and \abstractLabelRationale in the oracle setting (Row 3), and a one-point difference in the open setting (Row 6). The differences between \sentenceSelectionOnly and \sentenceSelectionLabel are larger, caused by examples where the model finds the evidence but fails to predict its relationship to the claim. We examine these in \S\ref{sec:error_analysis}.



\begin{table*}[t]
  \setlength{\tabcolsep}{2pt}
  \footnotesize
  \centering
  \begin{tabular}{L{0.12\linewidth} L{0.33\linewidth} L{0.44\linewidth} R{0.06\linewidth} }
    \toprule
    \textbf{Mistake type} &  \textbf{Claim} &  \textbf{Gold evidence} & \textbf{Count} \\
    \thinmid

    Science background & Rapamycin accelerates aging in fruit flies. & \refutesCovid{\dots feeding rapamycin to adult Drosophila produces life span extension \dots} & 29 / 37  \\
    \hdashline[0.4pt/2pt]

    Directionality & Inhibiting glucose-6-phospate dehydrogenase impairs lipogenesis & \supportsCovid{\dots suppression of 6PGD decreased lipogenesis \dots} & 9 / 37 \\

    \hdashline[0.4pt/2pt]

    Numerical reasoning & Bariatric surgery reduces resolution of diabetes. & \refutesCovid{Strong associations were found \dots with a HR of 9.29 (95\% CI 6.84-12.62)...} & 9 / 37 \\

    \hdashline[0.4pt/2pt]

    Cause and effect & Major vault protein (MVP) functions to decrease tumor aggression. & \refutesCovid{Knockout of MVP leads to miR-193a accumulation...inhibiting tumor progression} & 7 / 37  \\

    \hdashline[0.4pt/2pt]

    Coreference & Low saturated fat diets have adverse effects on the development of infants & \refutesCovid{Neurological development of children in the intervention group was at least as good as ... the control group} & 6 / 37 \\

    \arrayrulecolor{black}
    \bottomrule
  \end{tabular}
  \caption{Types of reasoning required to correct 37 mistaken predictions.}
  \label{tbl:error_analysis}

\end{table*} 
\subsection{Verifying claims about COVID-19} \label{sec:covid}
We conduct exploratory experiments using our system to verify claims concerning COVID-19. We tasked a medical student to write 36 COVID-related claims. For each claim $c$, we used \sysname to predict evidence abstracts $\Ehat(c)$. The annotator examined each $(c, \Ehat(c))$ pair. A pair was labeled \emph{plausible} if $\Ehat(c)$ was nonempty, and at least half of the evidence abstracts in $\Ehat(c)$ were judged to have reasonable rationales and labels. For 23 / 36 claims, the response of \sysname was deemed plausible by our annotator, demonstrating that \sysname is able to successfully retrieve and classify evidence in many cases. Two examples are shown in Table \ref{tbl:covid_example}.






\subsection{Error analysis} \label{sec:error_analysis}
To better understand the errors made by \sysname, we conduct a manual analysis of 37 predictions where an evidence abstract was correctly retrieved, but where the model failed to identify any relevant rationales or predicted an incorrect label. We identify five modeling capabilities, and count the number of missed examples where each modeling capability is required to make a correct prediction (Table \ref{tbl:error_analysis}). Many examples require more than one of these capabilities.

\begin{description}[leftmargin=*,leftmargin=0pt]
  \item \textbf{Science background} includes knowledge of domain-specific lexical relationships -- e.g. Drosophila is a synonym for fruit fly.
  \item \textbf{Directionality} requires understanding increases or decreases in scientific quantities -- e.g. decreased (but not increased) lipogenesis implies impairment of the lipogenesis process.
  \item \textbf{Numerical reasoning} involves interpreting numerical or statistical statements, often reporting confidence intervals or p-values.
  \item \textbf{Cause and effect} requires reasoning about counterfactuals -- e.g. if knocking out MVP inhibits tumor progression, then the presence of MVP enables tumor growth and aggression.
  \item \textbf{Coreference} involves drawing conclusions using context stated outside of a rationale sentence -- e.g. the ``intervention group'' mentioned in Table \ref{tbl:error_analysis} was on a low saturated fat diet.
\end{description}




 \section{Related work}

\label{sec:related_work}

\noindent {\bf Fact checking and rationalized NLP models}
Fact-checking datasets include PolitiFact \cite{Vlachos2014FactCT}, Emergent \cite{ferreira-vlachos-2016-emergent}, LIAR \cite{Wang2017LiarLP}, SemEval 2017 Task 8 RumorEval \cite{derczynski-etal-2017-semeval}, Snopes \cite{Popat2017WhereTT}, CLEF-2018 CheckThat! \cite{nakov2018overview}, Verify \cite{baly-etal-2018-integrating}, Perspectrum \cite{Chen2019SeeingTF}, \fever \cite{Thorne2018FEVERAL}, and \snopes \cite{Hanselowski2019ARA}. \cite{Hanselowski2019ARA} provides a thorough review.  To our knowledge, there are no existing data sets for scientific claim verification. We refer to our task as ``claim verification'' rather than ``fact-checking'' to emphasize that our focus is to help researchers make sense of scientific findings, not to counter disinformation.

Fact-checking is one of a number of tasks where a model is required to justify a prediction via \emph{rationales} from the source document. The ERASER dataset \cite{DeYoung2019ERASERAB} provides a suite of benchmark datasets for evaluating rationalized NLP models. \ours can serve as an additional testbed for these techniques.










\vspace{.2cm} \noindent {\bf Related scientific NLP tasks}
\label{sec:related_scinlp}
The \ours task is closely related to three other scientific NLP tasks. 
The \emph{citation contextualization} task \cite{Cohan2015MatchingCT,Jaidka2017TheCS} is to identify all spans in a cited document that are relevant to a particular citation in a citing document. Expert annotators achieved only 21.7\% inter-annotator agreement on the BioMedSumm dataset \cite{BiomedSumm}. 

The \emph{evidence inference} task \cite{Lehman2019InferringWM, DeYoung2020EvidenceI2}, involves predicting the effect of a medical intervention on a specified outcome. This task can be viewed as verifying a specific subset of scientific claims.


\emph{Systematic reviews} are the ``gold standard'' for assessing the truth of scientific findings\footnote{\url{https://www.cochranelibrary.com/about/about-cochrane-reviews}}. Research on \emph{automated evidence synthesis} \cite{Marshall2019TowardSR, Beller2018MakingPW, Tsafnat2014SystematicRA, Marshall2017AutomatingBE} seeks to automate some steps of the systematic review process. Most work on evidence synthesis has focused on extracting relevant documents and phrases such as PICO snippets \cite{Nye2018ACW}. We hope that systems for claim verification will serve as components in future evidence synthesis frameworks.
 \section{Conclusion}

Claim verification allows us to trace the sources and measure the veracity of scientific claims. These abilities have emerged as particularly important in the context of the current pandemic, and the broader reproducibility crisis in science. In this article, we formalize the task of scientific claim verification, and release a dataset (\ours) and models (\sysname) to support work on this task.

Scientific claim verification poses a set of unique challenges, pushing the limits of neural models on complex language understanding and reasoning. Domain-adaptation techniques show promise, but our findings suggest that additional work is necessary to improve the performance of end-to-end systems. We also demonstrate how claim verification might work in practice, by applying our system to the real-world problem of verifying claims related to COVID-19. We hope that these resources encourage others to expand upon our work, and to further shed light on the broader and more challenging goal of scientific document understanding.
 


\begin{thebibliography}{36}
\expandafter\ifx\csname natexlab\endcsname\relax\def\natexlab#1{#1}\fi

\bibitem{baly-etal-2018-integrating}
Ramy "Baly, Mitra Mohtarami, James Glass, Llu{\'\i}s M{\`a}rquez, Alessandro
  Moschitti, and Preslav" Nakov. 2018.
\newblock Integrating stance detection and fact checking in a unified corpus.
\newblock In \emph{North American Association for Computational Linguistics
  (NAACL)}.

\bibitem{nakov2018overview}
Alberto Barr{\'o}n-Cede{\~n}o, Tamer Elsayed, Reem Suwaileh,
  Llu{\'i}s~M{\`a}rquez i~Villodre, Pepa Atanasova, Wajdi Zaghouani, Spas
  Kyuchukov, Giovanni Da~San Martino, and Preslav Nakov. 2018.
\newblock Overview of the clef-2018 checkthat! lab on automatic identification
  and verification of political claims. task 2: Factuality.
\newblock In \emph{CLEF}.

\bibitem{Beller2018MakingPW}
Elaine Beller, Justin Clark, Guy Tsafnat, Clive~Elliott Adams, Heinz Diehl,
  Hans Lund, Mourad Ouzzani, Kristina Thayer, James Thomas, Tari Turner, J.~S.
  Xia, Karen~A. Robinson, and Paul~P Glasziou. 2018.
\newblock Making progress with the automation of systematic reviews: principles
  of the international collaboration for the automation of systematic reviews
  (icasr).
\newblock \emph{Systematic Reviews}, 7.

\bibitem{Beltagy2019SciBERTPC}
Iz~Beltagy, Kyle Lo, and Arman Cohan.
\newblock {S}ci{BERT}: A pretrained language model for scientific text.
\newblock In \emph{Empirical Methods in Natural Language Processing (EMNLP)}.

\bibitem{Beltagy2020LongformerTL}
Iz~Beltagy, Matthew~E. Peters, and Arman Cohan. 2020.
\newblock Longformer: The long-document transformer.
\newblock \emph{ArXiv}, abs/2004.05150.

\bibitem{BergKirkpatrick2012AnEI}
Taylor Berg-Kirkpatrick, David Burkett, and Dan Klein. 2012.
\newblock An empirical investigation of statistical significance in nlp.
\newblock In \emph{Empirical Methods in Natural Language Processing (EMNLP)}.

\bibitem{Chen2019SeeingTF}
Sihao Chen, Daniel Khashabi, Wenpeng Yin, Chris Callison-Burch, and Dan Roth.
  2019.
\newblock Seeing things from a different angle: Discovering diverse
  perspectives about claims.
\newblock In \emph{North American Association for Computational Linguistics
  (NAACL)}.

\bibitem{Cohan2015MatchingCT}
Arman Cohan, Luca Soldaini, and Nazli Goharian. 2015.
\newblock Matching citation text and cited spans in biomedical literature: a
  search-oriented approach.
\newblock In \emph{North American Association for Computational Linguistics
  (NAACL)}.

\bibitem{BiomedSumm}
Kevin~Bretonnel Cohen, Hoa~Trang Dang, Anita de~Waard, Prabha Yadav, and Lucy
  Vanderwende. 2014.
\newblock Tac 2014 biomedical summarization track.
\newblock \url{https://tac.nist.gov/2014/BiomedSumm/}.

\bibitem{derczynski-etal-2017-semeval}
Leon Derczynski, Kalina Bontcheva, Maria Liakata, Rob Procter, Geraldine Wong
  Sak~Hoi, and Arkaitz Zubiaga. 2017.
\newblock \href {https://doi.org/10.18653/v1/S17-2006} {{S}em{E}val-2017 task
  8: {R}umour{E}val: Determining rumour veracity and support for rumours}.
\newblock In \emph{Proceedings of the 11th International Workshop on Semantic
  Evaluation (SemEval-2017)}.

\bibitem{Devlin2019BERTPO}
Jacob Devlin, Ming-Wei Chang, Kenton Lee, and Kristina Toutanova. 2019.
\newblock Bert: Pre-training of deep bidirectional transformers for language
  understanding.
\newblock In \emph{North American Association for Computational Linguistics
  (NAACL)}.

\bibitem{DeYoung2019ERASERAB}
Jay DeYoung, Sarthak Jain, Nazneen~Fatema Rajani, Eric Lehman, Caiming Xiong,
  Richard Socher, and Byron~C. Wallace. 2020{\natexlab{a}}.
\newblock Eraser: A benchmark to evaluate rationalized nlp models.
\newblock In \emph{ACL}.

\bibitem{DeYoung2020EvidenceI2}
Jay DeYoung, Eric Lehman, Ben Nye, Iain~James Marshall, and Byron~C. Wallace.
  2020{\natexlab{b}}.
\newblock Evidence inference 2.0: More data, better models.
\newblock \emph{ArXiv}, abs/2005.04177.

\bibitem{Dror2018TheHG}
Rotem Dror, Gili Baumer, Segev Shlomov, and Roi Reichart. 2018.
\newblock The hitchhiker's guide to testing statistical significance in natural
  language processing.
\newblock In \emph{Association for Computational Linguistics (ACL)}.

\bibitem{Efron1993AnIT}
Bradley Efron and Robert Tibshirani. 1993.
\newblock An introduction to the bootstrap.

\bibitem{ferreira-vlachos-2016-emergent}
William Ferreira and Andreas Vlachos. 2016.
\newblock \href {https://doi.org/10.18653/v1/N16-1138} {{E}mergent: a novel
  data-set for stance classification}.
\newblock In \emph{North American Association for Computational Linguistics
  (NAACL)}.

\bibitem{Gururangan2020DONTST}
Suchin Gururangan, Ana Marasovic, Swabha Swayamdipta, Kyle Lo, Iz~Beltagy, Doug
  Downey, and Noah~A. Smith. 2020.
\newblock Don’t stop pretraining: Adapt language models to domains and tasks.
\newblock \emph{Association for Computational Linguistics (ACL)}.

\bibitem{Hanselowski2019ARA}
Andreas Hanselowski, Christian Stab, Claudia Schulz, Zile Li, and Iryna
  Gurevych. 2019.
\newblock A richly annotated corpus for different tasks in automated
  fact-checking.
\newblock In \emph{Proceedings of the 23rd Conference on Computational Natural
  Language Learning (CoNLL)}.

\bibitem{Jaidka2017TheCS}
Kokil Jaidka, Muthu~Kumar Chandrasekaran, Devanshu Jain, and Min-Yen Kan. 2017.
\newblock The cl-scisumm shared task 2017: Results and key insights.
\newblock In \emph{BIRNDL@JCDL}.

\bibitem{Lehman2019InferringWM}
Eric Lehman, Jay DeYoung, Regina Barzilay, and Byron~C. Wallace. 2019.
\newblock Inferring which medical treatments work from reports of clinical
  trials.
\newblock In \emph{North American Association for Computational Linguistics
  (NAACL)}.

\bibitem{Lei2016RationalizingNP}
Tao Lei, Regina Barzilay, and Tommi~S. Jaakkola. 2016.
\newblock Rationalizing neural predictions.

\bibitem{Liu2019RoBERTaAR}
Yinhan Liu, Myle Ott, Naman Goyal, Jingfei Du, Mandar Joshi, Danqi Chen, Omer
  Levy, Mike Lewis, Luke Zettlemoyer, and Veselin Stoyanov. 2019.
\newblock Roberta: A robustly optimized bert pretraining approach.
\newblock \emph{ArXiv}, abs/1907.11692.

\bibitem{Lo2019GORCAL}
Kyle Lo, Lucy~Lu Wang, Mark Neumann, Rodney Kinney, and Daniel~S. Weld. 2020.
\newblock \href {https://arxiv.org/abs/1911.02782} {{S2ORC: The Semantic
  Scholar Open Research Corpus}}.
\newblock In \emph{Association for Computational Linguistics (ACL)}.

\bibitem{Marshall2017AutomatingBE}
Iain~James Marshall, Jo{\"e}l Kuiper, Edward Banner, and Byron~C. Wallace.
  2017.
\newblock Automating biomedical evidence synthesis: Robotreviewer.
\newblock \emph{Association for Computational Linguistics (ACL)}.

\bibitem{Marshall2019TowardSR}
Iain~James Marshall and Byron~C. Wallace. 2019.
\newblock Toward systematic review automation: a practical guide to using
  machine learning tools in research synthesis.
\newblock \emph{Systematic Reviews}, 8.

\bibitem{Nakov2004CitancesC}
Preslav~I Nakov, Ariel~S Schwartz, and Marti Hearst. 2004.
\newblock Citances: Citation sentences for semantic analysis of bioscience
  text.
\newblock In \emph{Proceedings of the SIGIR’04 workshop on Search and
  Discovery in Bioinformatics}.

\bibitem{Nye2018ACW}
Benjamin Nye, Junyi~Jessy Li, Roma Patel, Yinfei Yang, Iain~James Marshall, Ani
  Nenkova, and Byron~C. Wallace. 2018.
\newblock A corpus with multi-level annotations of patients, interventions and
  outcomes to support language processing for medical literature.
\newblock \emph{Proceedings of the conference. Association for Computational
  Linguistics. Meeting}, 2018:197--207.

\bibitem{Popat2017WhereTT}
Kashyap Popat, Subhabrata Mukherjee, Jannik Str{\"o}tgen, and Gerhard Weikum.
  2017.
\newblock Where the truth lies: Explaining the credibility of emerging claims
  on the web and social media.
\newblock In \emph{Proceedings of the International World Wide Web Conference
  (WWW)}.

\bibitem{Schuster2019TowardsDF}
Tal Schuster, Darsh~J. Shah, Yun Jie~Serene Yeo, Daniel Filizzola, Enrico
  Santus, and Regina Barzilay. 2019.
\newblock Towards debiasing fact verification models.
\newblock In \emph{Empirical Methods in Natural Language Processing (EMNLP)}.

\bibitem{Soleimani2019BERTFE}
Amir Soleimani, Christof Monz, and Marcel Worring. 2019.
\newblock Bert for evidence retrieval and claim verification.
\newblock \emph{ArXiv}, abs/1910.02655.

\bibitem{Thorne2018FEVERAL}
James Thorne, Andreas Vlachos, Christos Christodoulopoulos, and Arpit Mittal.
  2018.
\newblock Fever: a large-scale dataset for fact extraction and verification.
\newblock In \emph{North American Association for Computational Linguistics
  (NAACL)}.

\bibitem{Tsafnat2014SystematicRA}
Guy Tsafnat, Paul~P Glasziou, Miew~Keen Choong, Adam~G. Dunn, Filippo Galgani,
  and Enrico~W. Coiera. 2014.
\newblock Systematic review automation technologies.
\newblock \emph{Systematic Reviews}, 3:74 -- 74.

\bibitem{Vlachos2014FactCT}
Andreas Vlachos and Sebastian Riedel. 2014.
\newblock \href {https://doi.org/10.3115/v1/W14-2508} {Fact checking: Task
  definition and dataset construction}.
\newblock In \emph{Proceedings of the {ACL} 2014 Workshop on Language
  Technologies and Computational Social Science}.

\bibitem{Wang2020CORD19TC}
Lucy~Lu Wang, Kyle Lo, Yoganand Chandrasekhar, Russell Reas, Jiangjiang Yang,
  Darrin Eide, Kathryn Funk, Rodney Kinney, Ziyang Liu, William. Merrill, Paul
  Mooney, Dewey~A. Murdick, Devvret Rishi, Jerry Sheehan, Zhihong Shen, Brandon
  Stilson, Alex~D. Wade, Kuansan Wang, Christopher Wilhelm, Boya Xie,
  Douglas~M. Raymond, Daniel~S. Weld, Oren Etzioni, and Sebastian Kohlmeier.
  2020.
\newblock Cord-19: The covid-19 open research dataset.
\newblock \emph{ArXiv}, abs/2004.10706.

\bibitem{Wang2017LiarLP}
William~Yang Wang. 2017.
\newblock \href {https://doi.org/10.18653/v1/P17-2067} {{``}liar, liar pants on
  fire{''}: A new benchmark dataset for fake news detection}.
\newblock In \emph{Association for Computational Linguistics (ACL)}.
  Association for Computational Linguistics.

\bibitem{Wolf2019HuggingFacesTS}
Thomas Wolf, Lysandre Debut, Victor Sanh, Julien Chaumond, Clement Delangue,
  Anthony Moi, Pierric Cistac, Tim Rault, R'emi Louf, Morgan Funtowicz, and
  Jamie Brew. 2019.
\newblock Huggingface's transformers: State-of-the-art natural language
  processing.
\newblock \emph{ArXiv}, abs/1910.03771.

\end{thebibliography}
 
\bibliographystyle{acl_natbib}

\clearpage

\appendix

\section{Model implementation details} \label{sec:model_details}

All models are implemented using the Huggingface Transformers package \cite{Wolf2019HuggingFacesTS}.



\subsection{Parameters for the final \sysname system} \label{sec:parameters}

For the \componentone module, \sysname retrieves the top $k=3$ documents ranked by TF-IDF similarity using unigram + bigram features. These parameters are tuned on the \ours development set.

When making predictions using the \componenttwo module described in \S\ref{sec:rationale_selection}, we find that the usual decision rule of predicting $\hat{z}_i = 1$ when $\tilde{z}_i \geq 0.5$ works well for models trained on \ours. However, for models trained on \fever and \snopes, we achieve better performance by tuning the classification threshold $t$, such that $\hat{z}_i = 1$  when $\tilde{z}_i \geq t$, on the \ours dev set. The best threshold was $t=0.025$ when training on \fever, and $t=0.75$ when training on \snopes.





\subsection{Training the \componenttwo module}

We experiment with various learning rates when training \scibert, \bioroberta, \robertabase, and \robertalarge. Below we describe the setting for training \robertalarge.

For models trained on \ours, we use an initial learning rate of 1e-5 on the transformer base and 1e-3 on the linear layer. For \fever + \ours, the learning rate is set to 1e-5 for the entire model for pre-training on \fever and fine-tuning on \ours. We use a batch size of 256 through gradient accumulation and apply cosine learning rate decay over 20 epochs to find the best performing model on the dev set.

For models trained on \fever, we set the learning rate to 5e-6 for the transformer base and 5e-5 for the linear layer. For models trained on \snopes, we set the learning rate 1e-5 for the transformer base and 1e-4 for the linear layer. We find that these learning rates help the models converge. We only train the model for 3 epochs on \fever and 5 epochs on \snopes because they are larger datasets and the models converged within early epochs.



\subsection{Training the \componentthree module}

We adopt similar settings as we used for the \componenttwo module and only change the learning rate to 1e-5 for the transformer base and 1e-4 for the linear layer for models trained on \ours, \fever, and \snopes. When training on claim / cited abstract pairs labeled \notenough, we use the $k$ sentences in the abstract with greatest similarity to the claim as rationales (\S\ref{sec:baselines}). $k$ is sampled from $\{0, 1\}$ with uniform probability.




\subsection{Additional Training details}

All models are trained using a single Nvidia P100 GPU on Google Colabortoary Pro platform.\footnote{\url{https://colab.research.google.com/}} For the \componenttwo module, it takes about 150 minutes to train on \ours for 20 epochs. 120 minutes on \snopes for 5 epochs, and 700 minutes on \fever for 3 epochs. For the \componentthree module, it takes about 130 minutes to train on \ours for 20 epochs, 160 minutes on \snopes for 5 epochs, and 640 minutes on \fever for 3 epochs.




\subsection{Hyperparameters search}

The learning rate, batch size, and number of epochs are the most important hyperparameters. We perform manual tuning and select the hyperparameters that produce the highest F1 on the development set. For the learning rate, we experiment with 1e-3, 1e-4, 5e-5, 1e-5, and 5e-6. For batch size, we experiment with 64 and 256. The number of epochs are cutoff after the model converges.



\section{Full task: additional performance measurements} \label{sec:uncertainty}

We assess the uncertainty in the results reported in the main results (Table \ref{tbl:main_results}) using a simple bootstrap approach \cite{Dror2018TheHG, BergKirkpatrick2012AnEI, Efron1993AnIT}. Given our test set with $n_{\textrm{test}} = 300$ claims, we generate $n_{\textrm{boot}} = 10,000$ bootstrap-resampled test sets by resampling (uniformly, with replacement) $n_{\textrm{test}}$ claims from the test set. For each resampled test set, we compute the metrics in Table \ref{tbl:main_results}. Table \ref{tbl:results_test_bootstrap} reports the mean and standard deviation of these metrics, computed over the bootstrap samples. Table \ref{tbl:results_dev_bootstrap} reports dev set metrics. Our conclusion that training on \ours improves performance is robust to the uncertainties presented in these tables.

\clearpage

\begin{table*}[t]
  \footnotesize
  \setlength{\tabcolsep}{0.5em}


  \begin{subtable}[h]{\linewidth}
    \centering
    \begin{tabularx}{0.8\linewidth}{*{3}{l} *{6}{c}}

      \toprule

      & & & \multicolumn{6}{c}{\textbf{Sentence-level}} \\
      & & & \multicolumn{3}{c}{\textbf{Selection-Only}} & \multicolumn{3}{c}{\textbf{Selection+Label}} \\

      \cmidrule(lr){4-6} \cmidrule(lr){7-9}

      Retrieval & Model & Row & P & R & F1 & P & R & F1 \\
      \midrule
      \multirow{4}{*}{\shortstack[c]{\textbf{Oracle} \\ \textbf{abstract}}} &
      Oracle rationale & 1 & $100.0_{0.0}$ & $80.5_{3.3}$ & $89.2_{2.1}$ & $89.6_{2.7}$ & $72.2_{3.7}$ & $79.9_{3.0}$ \\

      \arrayrulecolor{black!30}\cmidrule(lr){2-9}
      & Zero-shot & 2 & $42.6_{2.2}$ & $45.2_{3.2}$ & $43.8_{2.0}$ & $36.2_{2.5}$ & $38.4_{3.0}$ & $37.2_{2.3}$ \\
      & \sysname & 3 & $76.2_{2.9}$ & $63.9_{3.6}$ & $69.4_{2.6}$ & $66.5_{3.4}$ & $55.7_{3.7}$ & $60.6_{3.1}$ \\

      \thickmid

      \multirow{4}{*}{\textbf{Open}} &
      Oracle rationale               & 4 & $100.0_{0.0}$ & $56.6_{4.0}$ & $72.2_{3.3}$ & $87.6_{3.5}$ & $49.5_{3.9}$ & $63.2_{3.7}$ \\
      \arrayrulecolor{black!30}\cmidrule(lr){2-9}
      & Zero-shot & 5 & $28.7_{2.3}$ & $37.6_{3.4}$ & $32.5_{2.3}$ & $23.8_{2.3}$ & $31.1_{3.1}$ & $26.9_{2.3}$ \\
      & \sysname & 6 & $45.0_{3.0}$ & $47.4_{3.8}$ & $46.1_{3.0}$ & $38.5_{3.0}$ & $40.6_{3.6}$ & $39.5_{3.0}$ \\

    \end{tabularx}
    \subcaption{Sentence-level results.}
  \end{subtable}

  \begin{subtable}[h]{\linewidth}
    \centering
    \begin{tabularx}{0.8\linewidth}{*{3}{l} *{6}{c}}

      \toprule

      &  &  & \multicolumn{6}{c}{\textbf{Abstract-level}} \\
      &  &  & \multicolumn{3}{c}{\textbf{Label-Only}} & \multicolumn{3}{c}{\textbf{Label+Rationale}} \\

      \cmidrule(lr){4-6} \cmidrule(lr){7-9}

      Retrieval & Model & Row & P & R & F1 & P & R & F1 \\

      \midrule

      \multirow{4}{*}{\shortstack[c]{\textbf{Oracle} \\ \textbf{abstract}}} &
      Oracle rationale & 1 & $90.1_{2.2}$ & $77.5_{2.8}$ & $83.3_{2.4}$ & $90.1_{2.2}$ & $77.5_{2.8}$ & $83.3_{2.4}$ \\

      \arrayrulecolor{black!30}\cmidrule(lr){2-9}

      & Zero-shot & 2 & $86.9_{2.9}$ & $53.6_{3.4}$ & $66.3_{3.1}$ & $67.9_{3.9}$ & $41.9_{3.2}$ & $51.8_{3.4}$ \\
      & \sysname & 3 & $87.3_{2.6}$ & $65.3_{3.2}$ & $74.7_{2.8}$ & $84.9_{2.8}$ & $63.5_{3.2}$ & $72.6_{2.9}$ \\

      \thickmid

      \multirow{4}{*}{\textbf{Open}} &

      Oracle rationale               & 4 & $88.9_{2.7}$ & $54.1_{3.5}$ & $67.2_{3.2}$ & $88.9_{2.7}$ & $54.1_{3.5}$ & $67.2_{3.2}$ \\

      \arrayrulecolor{black!30}\cmidrule(lr){2-9}
      & Zero-shot & 5 & $56.0_{3.9}$ & $42.3_{3.4}$ & $48.2_{3.3}$ & $42.3_{4.0}$ & $32.0_{3.2}$ & $36.4_{3.3}$ \\
      & \sysname & 6  &  $47.5_{3.3}$ & $47.3_{3.5}$ & $47.4_{3.1}$ & $46.6_{3.3}$ & $46.4_{3.5}$ & $46.4_{3.1}$ \\

      \arrayrulecolor{black}\bottomrule

    \end{tabularx}
    \subcaption{Abstract-level results}
  \end{subtable}

  \caption{Test set results as in Table \ref{tbl:main_results}, reporting mean and standard deviation over 10,000 bootstrap samples. Standard deviations are reported as subscripts. Some means reported here are slightly different from Table \ref{tbl:main_results} due to sampling variability.}

  \label{tbl:results_test_bootstrap}


\end{table*} \begin{table*}[h!]
  \footnotesize
  \setlength{\tabcolsep}{0.5em}


  \begin{subtable}[h]{\linewidth}
    \centering
    \begin{tabularx}{0.8\linewidth}{*{3}{l} *{6}{c}}

      \toprule

      & & & \multicolumn{6}{c}{\textbf{Sentence-level}} \\
      & & & \multicolumn{3}{c}{\textbf{Selection-Only}} & \multicolumn{3}{c}{\textbf{Selection+Label}} \\

      \cmidrule(lr){4-6} \cmidrule(lr){7-9}

      Retrieval & Model & Row & P & R & F1 & P & R & F1 \\
      \midrule
      \multirow{4}{*}{\shortstack[c]{\textbf{Oracle} \\ \textbf{abstract}}} &
      Oracle rationale & 1 &   $100.0_{0.0}$ &       $81.9_{3.2}$ &       $90.0_{1.9}$ &   $91.4_{2.5}$ &   $74.9_{3.6}$ &   $82.3_{2.9}$ \\

      \arrayrulecolor{black!30}\cmidrule(lr){2-9}
      & Zero-shot & 2 &  $40.7_{2.1}$ &       $48.1_{3.4}$ &       $44.0_{2.1}$ &   $36.1_{2.5}$ &   $42.6_{3.4}$ &   $39.0_{2.5}$ \\
      & \sysname & 3 &  $79.4_{2.7}$ &       $59.0_{3.6}$ &       $67.7_{2.8}$ &   $71.4_{3.5}$ &   $53.0_{3.6}$ &   $60.8_{3.3}$ \\

      \thickmid

      \multirow{4}{*}{\textbf{Open}} &
      Oracle rationale               & 4 &   $100.0_{0.0}$ &       $58.4_{4.3}$ &       $73.7_{3.4}$ &   $90.2_{3.3}$ &   $52.7_{4.3}$ &   $66.4_{3.9}$ \\
      \arrayrulecolor{black!30}\cmidrule(lr){2-9}
      & Zero-shot & 5 &   $28.6_{2.0}$ &       $38.5_{3.6}$ &       $32.8_{2.3}$ &   $24.8_{2.2}$ &   $33.4_{3.4}$ &   $28.4_{2.4}$ \\
      & \sysname & 6 &       $52.5_{3.5}$ &       $43.8_{3.7}$ &       $47.7_{3.2}$ &   $46.9_{3.7}$ &   $39.2_{3.6}$ &   $42.6_{3.2}$ \\

    \end{tabularx}
    \subcaption{Sentence-level results.}
  \end{subtable}

  \begin{subtable}[h]{\linewidth}
    \centering
    \begin{tabularx}{0.8\linewidth}{*{3}{l} *{6}{c}}

      \toprule

      &  &  & \multicolumn{6}{c}{\textbf{Abstract-level}} \\
      &  &  & \multicolumn{3}{c}{\textbf{Label-Only}} & \multicolumn{3}{c}{\textbf{Label+Rationale}} \\

      \cmidrule(lr){4-6} \cmidrule(lr){7-9}

      Retrieval & Model & Row & P & R & F1 & P & R & F1 \\

      \midrule

      \multirow{4}{*}{\shortstack[c]{\textbf{Oracle} \\ \textbf{abstract}}} &
      Oracle rationale & 1 &  $91.4_{2.2}$ &   $76.1_{3.0}$ &   $83.0_{2.5}$ &       $91.4_{2.2}$ &       $76.1_{3.0}$ &       $83.0_{2.5}$ \\

      \arrayrulecolor{black!30}\cmidrule(lr){2-9}

      & Zero-shot & 2 &    $88.9_{2.8}$ &   $58.3_{3.7}$ &   $70.4_{3.2}$ &       $69.2_{3.9}$ &       $45.4_{3.5}$ &       $54.8_{3.5}$ \\
      & \sysname & 3 &    $91.0_{2.3}$ &   $67.4_{3.3}$ &   $77.4_{2.7}$ &       $85.2_{2.9}$ &       $63.2_{3.5}$ &       $72.5_{3.1}$ \\

      \thickmid

      \multirow{4}{*}{\textbf{Open}} &

      Oracle rationale               & 4 &    $91.0_{2.6}$ &   $53.1_{3.8}$ &   $67.0_{3.4}$ &       $91.0_{2.6}$ &       $53.1_{3.8}$ &       $67.0_{3.4}$ \\

      \arrayrulecolor{black!30}\cmidrule(lr){2-9}
      & Zero-shot & 5 &   $52.7_{3.7}$ &   $41.6_{3.7}$ &   $46.5_{3.4}$ &       $43.6_{3.7}$ &       $34.4_{3.5}$ &       $38.4_{3.3}$ \\
      & \sysname & 6  &   $55.4_{3.7}$ &   $47.5_{3.6}$ &   $51.0_{3.3}$ &       $52.6_{3.7}$ &       $45.1_{3.6}$ &       $48.5_{3.3}$ \\

      \arrayrulecolor{black}\bottomrule

    \end{tabularx}
    \subcaption{Abstract-level results}
  \end{subtable}

  \caption{Dev set results as in Table \ref{tbl:main_results}, reporting mean and standard deviation over 10,000 bootstrap samples.}

  \label{tbl:results_dev_bootstrap}

\end{table*}
 
\clearpage




\section{Dataset collection and corpus statistics} \label{sec:data_collection}


\subsection{Annotation examples}

\paragraph{Converting citances to claims} Figure \ref{fig:claim} shows an example of a citance re-written as a claim. The citance discusses the relationship between ``atherosclerotic CVD'' and ``IL-6'', and cites two papers (\textbf{44} and \textbf{45}) as evidence. To convert to a claim, the acronym ``CVD'' is expanded to ``cardiovascular disease'', irrelevant information is removed, and the claim is written as an atomic factual statement.

\begin{figure}[t]
  \centering
  \includegraphics[width=\columnwidth]{figures/claim-fig.pdf}

  \caption{A claim written based on a citance. Material unrelated to the citation is removed. The acronym ``CVD'' is expanded to ``cardiovascular disease''.}
  \label{fig:claim}
\end{figure} 
\paragraph{Multiple rationales} Figure \ref{fig:multiple_rationales} shows a claim supported by two rationales from the same abstract. The text of each rationale on its own is sufficient to entail the claim.

\begin{figure}[t]
  \centering

  \includegraphics[width=\columnwidth]{figures/multiple-rationales-fig.pdf}
  \caption{A claim supported by two rationales from the same abstract. The text of each rationale on its own provides sufficient evidence to verify the claim.}
  \label{fig:multiple_rationales}

\end{figure} 



\subsection{Annotators and quality control}

\paragraph{Claim writing} The claim writers included four experts with background in scientific NLP, fifteen undergraduates studying the life sciences, and four graduate students (doctoral or medical) in the life sciences. Student claim writers attended an in-person training session where they were introduced to the task and received in-person feedback from the four experts. Following training, student annotators continued writing claims remotely. The expert annotators monitored claims for quality during the remote annotation process, and provided feedback when necessary; low-quality claims were returned to the annotators for re-writing. As a final check, all submitted claims were proofread (and edited if necessary) by an undergraduate whose claims were deemed especially high-quality by the expert annotators.

\paragraph{Claim negations}

As mentioned in \S\ref{sec:claim_writing}, an expert annotator wrote claim negations to introduce cases where an abstract \refutes a claim. The annotator skipped claims that could only be negated by adding obvious triggers like ``not''. The majority of claim negations involved a reversal of effect direction; for instance ``\emph{A high microerythrocyte count protects against severe anemia}'' can be negated as ``\emph{A high microerythrocyte count raises vulnerability to severe anemia}''.

\paragraph{Claim verification}

The verifiers included three NLP experts, five undergraduates studying life sciences, and five graduate students studying life sciences. Annotations were performed remotely through a web interface. Annotators were required to pass a 10-question ``quiz'' before annotating their own claims. After passing the quiz, subsequent submissions were reviewed by an NLP expert until that expert deemed the annotator reliable. Approved annotators were then assigned to review each others' submissions. In general, graduate students were assigned to review annotations from undergraduates.





\subsection{Corpus} \label{sec:data_source_appendix}

\paragraph{Additional statistics} Table \ref{tbl:journal_counts} shows the number of cited abstracts from each of our selected journals. The ``Other'' category includes ``co-cited'' (\S\ref{sec:data_source}) abstracts that came from journals not among our pre-defined set.

\begin{table}[t]
  \centering

  \begin{tabular}{l r}
  \toprule
  Journal & Count \\
  \thickmid
  BMJ                            &     60 \\
Blood                          &      8 \\
Cancer Cell                    &      8 \\
Cell                           &     51 \\
Cell Metabolism                &     10 \\
Cell Stem Cell                 &     41 \\
Circulation                    &     12 \\
Immunity                       &     33 \\
JAMA                           &     79 \\
Molecular Cell                 &     27 \\
Molecular Systems Biology      &      5 \\
Nature                         &     29 \\
Nature Cell Biology            &     26 \\
Nature Communications          &     19 \\
Nature Genetics                &      8 \\
Nature Medicine                &     89 \\
Nature Methods                 &      1 \\
Nucleic Acids Research         &     10 \\
Plos Biology                   &     36 \\
Plos Medicine                  &     38 \\
Science                        &      7 \\
Science Translational Medicine &      2 \\
The Lancet                     &     22 \\
\thinmid
Other                          &    120 \\

\thickmid

Total                          &    741 \\
\bottomrule
\end{tabular}

  \caption{Number of cited documents by journal. Some co-cited articles (\S \ref{sec:data_source}) come from journals outside our curated set; these are indicated by ``Other''.}
  \label{tbl:journal_counts}
\end{table} 
Figure \ref{fig:mesh_terms} shows counts for the most-frequent MESH terms\footnote{\url{https://www.ncbi.nlm.nih.gov/mesh}} among the cited abstracts. These terms reflect studies ranging from clinical trial reports (``Humans'', ``Middle Aged'', ``Child'') to molecular biology (``Cell Line'', ``RNA'', ``Mutation'').

\begin{figure}[t]
  \centering
  \includegraphics[width=\columnwidth]{figures/mesh-terms-fig.pdf}

  \caption{Fraction of evidence abstracts in which each MESH term occurs.}
  \label{fig:mesh_terms}
\end{figure} 
\paragraph{Distractor abstracts} In \S\ref{sec:data_source}, we mention how we increase the size of the corpus by adding \emph{distractor abstracts}. The reason why we do not use the entirety of a large research corpus like S2ORC as our fact-checking corpus is that doing so would introduce many \emph{false negative retrievals}: abstracts containing evidence relevant to a given claim, but not mentioned in the claim's source citance. This can occur either because the citance authors simply were not aware of these abstracts, or because the abstracts were published after the citance was written. These retrievals would be incorrectly marked wrong in by evaluation metrics.

Distractor abstracts as defined in \S\ref{sec:data_source} have two qualities that make them a good addition to the \ours corpus: (1) They are cited in the same articles as our evidence abstracts, meaning that they often discuss similar topics and increase the difficulty of abstract retrieval methods based on lexical similarity. (2) The authors of our citances were aware of the distractor abstracts, and chose not to mention them in the citances used to generate claims. This makes them unlikely to be a source of false positive retrievals.




\section{Annotation interfaces and guidelines} \label{sec:annotation_interfaces}

We show a screenshot of the claim interface in Figure \ref{fig:claim_interface}, and evidence interface in Figure \ref{fig:evidence_interface}.

\begin{figure*}[t]
  \centering
  \includegraphics[width=\textwidth]{figures/claim-interface-fig.pdf}

  \caption{The claim-writing interface. The citation sentence is highlighted in blue on the top left. Additional context is provided on bottom left. The right side shows two claims that could be written based on this citation sentence.}
  \label{fig:claim_interface}
\end{figure*} 
\begin{figure*}[h]
  \centering
  \includegraphics[width=\textwidth]{figures/evidence-interface-fig.pdf}

  \caption{The evidence collection interface.}
  \label{fig:evidence_interface}
\end{figure*}  

\end{document}